\documentclass{article}
\usepackage{../assignment_preamble}

\title{Homework 3}
\author{Ravi Kini}
\date{October 15, 2025}

\begin{document}

\maketitle

\problem{}
\begin{center}
\begin{tikzpicture}[
       decoration = {markings,
                     mark=at position .5 with {\arrow{Stealth[length=2mm]}}},
       dot/.style = {circle, fill, inner sep=2.4pt, node contents={},
                     label=#1},
every edge/.style = {draw, postaction=decorate}
                        ]
\node (a) at (1.5,0) [dot=above:$A$];
\node (b) at (-1.5,0) [dot=above:$B$];
\node (c) at (0,-1) [dot=right:$C$];
\node (d) at (2,-1) [dot=right:$D$];
\node (h) at (-2,-1) [dot=left:$H$];
\node (e) at (-1.5,-2) [dot=below:$E$];
\node (f) at (0,-2) [dot=below:$F$];
\node (g) at (1.5,-2) [dot=below:$G$];
%
\path
(a) edge (h)    (h) edge (f)
(a) edge (c)    (c) edge (e)
(a) edge (c)    (c) edge (f)     
(b) edge (c)    (c) edge (f) 
(b) edge (d)    (d) edge (f)  
(b) edge (d)    (d) edge (g);
\end{tikzpicture}
\end{center}
\begin{align*}
    A & = \{0, 1\} \to 0 \\
    B & = \{0, 1\} \to 1 \\
    C & = \max(A, B) \to 1 \ \\
    D & = B \to 1 \\
    E & = C \to 1 \\
    F & = \min(H, 1 - \max(C, D)) \to 0 \\
    G & = D \to 1 \\
    H & = 1 - A \to 1 \\
\end{align*}
\subproblem{(a)}
Changing $B$ potentially affects the values of $C$, $D$, $F$, and $G$. If $B = 0$, $C = \max(0, 0) = 0$ and $D = 0$. Consequently, $F = \min(1, 1 - \max(0, 0)) = 1$.
\subproblem{(b)}
Changing $B$ and $D$ potentially affects the values of $C$, $F$, and $G$. If $B = 0$ and $D = 1$, $C = \max(0, 0) = 0$. Consequently, $F = \min(1, 1 - \max(0, 1)) = 0$.
\clearpage

\problem{}
\subproblem{(a)}
If your PhD application is a good fit for the program you apply to, you are highly likely to get admitted. However, if your application is a good fit and there are several other applicants with similar profiles, it is not highly likely that you will get admitted.

\subproblem{(b)}
Let $X$ represent the quality of your PhD application ($0$ if a bad fit, $1$ if a good fit) and $Z$ represent the amount of strong competition you face ($0$ if low, $1$ if high). Letting $Y$ represent the likelihood of you being admitted ($0$ if low, $1$ if high), this can be modelled as $Y = \min(X, 1 - Z)$. In the original causal world, let $Z = 0$. If $X = 1$, $Y = \min(1, 1 - 0) = 1$. However, if $X = 1$ and $Z = 1$, $Y = \min(1, 1 - 1) = 0$.

\clearpage

\problem{}
If $X = 0$, then $Y = 0$ and $Z = \min(0, 1 - 0) = 0$. If $X = 0$ and $Y = 0$, then $Z = \min(0, 1 - 0) = 0$. If $X = 1$, then $Y = 0$ and $Z = \min(1, 1 - 0) = 1$. If $X = 1$ and $Y = 0$, then $Z = \min(1, 1 - 0) = 0$. In both cases, the conclusion holds based on the premises. Note that it is impossible for $X = 1$ and $Y = 1$ or $X = 0$ and $Y = 1$, as that would violate the second premise. As it is always th case that the conclusion is true or one of the premises is false, the causal world fails to be a formal counterexample to cautious monotonicity.

\clearpage

\problem{}
\begin{center}
\begin{tikzpicture}[
       decoration = {markings,
                     mark=at position .5 with {\arrow{Stealth[length=2mm]}}},
       dot/.style = {circle, fill, inner sep=2.4pt, node contents={},
                     label=#1},
every edge/.style = {draw, postaction=decorate}
                        ]
\node (d) at (0,0) [dot=above:$D$];
\node (a) at (-1,-1) [dot=left:$A$];
\node (b) at (1,-1) [dot=right:$B$];
\node (e) at (0,-2) [dot=below:$E$];
%
\path
(d) edge (a)    (a) edge (e)
(d) edge (b)    (b) edge (e);
\end{tikzpicture}
\end{center}

$D$ takes on values in the range $[0, 1]$. Under the top left mapping, with weight $-34.4$ and bias $2.14$, this is mapped to a range of $[-32.26, 2.14]$. Under the softplus function, the top variable ($A$) takes on values in the range $[9.780 \cdot 10^{-15}, 2.251]$. Under the top right mapping, with weight $-1.30$, this is mapped to a range of $[-2.927, -1.270 \cdot 10^{-14}]$. Under the bottom left mapping, with weight $-2.52$ and bias $1.29$, this is mapped to a range of $[-1.23, 1.29]$. Under the softplus function, the bottom variable ($B$) taken on values in the range $[0.256, 1.533]$. Under the bottom right mapping, with weight $2.28$, this is mapped to a range of $[0.585, 3.496]$. Summing and adding the final bias of $-0.58$, $E$ takes on values in the range $[-2.922, 2.916]$.

\begin{align*}
    D & \in [0, 1] \\
    A & = \ln(1 + e^{-34.4 \cdot A + 2.14}) \\
    B & = \ln(1 + e^{-2.52 \cdot A + 1.29}) \\
    E & = -1.30 \cdot A + 2.28 \cdot B + -0.58 \\
\end{align*}
\end{document}